% !TeX root = ../main.tex
\section{Working Without the Dry Signal}
\label{sec:nodi}

\subsection{Unpaired and blind modelling}

Wright, Välimäki and Juvela posed amplifier modelling as domain transfer from
unprocessed to processed guitar and trained a real-time WaveNet only
adversarially, against spectral or MelGAN discriminators~\cite{wright2023adversarial}.
With input and target from the same guitar, the best unpaired model scored
83--92 MUSHRA across three plugin tones; with a different guitar for the
input, the realistic case, it scored 77--84, comparable to an oracle baseline
that equalised the true plugin's output (72--81). Training took 4--10 GPU
hours per tone. Chen et al.\ improved stability with multi-scale and
multi-period discriminators and extra unaligned clean data, and found
loudness normalisation essential~\cite{chen2024unsupervised}.

Moliner et al.\ estimate the clean signals and the effect jointly by diffusion
posterior sampling with a prior trained on clean guitar, and compare this
with adversarial training for black-box (GCN, S4) and grey-box
Wiener--Hammerstein operators~\cite{moliner2025unsupervised}. Diffusion stays
stable with only 18\,s of processed audio, where adversarial training
degrades; adversarial training is better on heavy distortion; and the grey-box
operator matched the black-box ones under diffusion but collapsed under
adversarial training. Švento et al.\ estimate a memoryless nonlinearity
(a spline) together with the clean signal~\cite{svento2025estimation}, which
gives both a dry estimate and a waveshaper curve, and LOUDAR moves the
operator into an autoencoder's latent space~\cite{svento2026loudar}.

All of these were evaluated on isolated, plugin-rendered guitar, mostly with
distortion alone. None handles cabinet, modulation, delay and reverberation
together, and none was tested on a separated stem of a commercial mix.

\subsection{Effect removal}

A recovered dry signal turns the wet-only problem back into a paired one.
Imort et al.\ removed distortion with separation networks~\cite{imort2022distortion}.
RemFX detects which of five effects are present and chains effect-specific
removal models; SI-SDR falls from 16.7\,dB with one effect to 2.6\,dB with
five~\cite{rice2023general}. Distortion Recovery restores a mel spectrogram and
re-synthesises it with a vocoder, trained on commercial amplifier presets and
4.5\,h of real amplifier recordings~\cite{lee2024distortion}. EG-VAE factorises
a processed recording into frame-level content and a global tone embedding;
swapping the embedding transfers the tone, and a learned masking mode removes
it~\cite{yeh2026egvae}. On plugin tones it reaches a tone-similarity opinion
score of 4.30 on seen and 3.70 on unseen tones (oracle 4.48; baselines at most
2.74), and a ``dryness'' score of 3.42--3.50 for removal against 4.39--4.58
for the true DI. A baseline that fitted a fixed differentiable chain of generic
blocks to EG-VAE's dry estimate scored only 1.44--2.74 in tone similarity,
evidence that generic equaliser, compressor and distortion blocks do not cover
guitar tone diversity. No code or weights were found for EG-VAE.

Removal and remodelling have been chained. Hinrichs et al.\ remove a single
effect with a separation network, then train a neural model on the recovered
pair; listeners rated the copies at medians of 92--100 for slapback, phaser
and soft clipping, but larger models lost the effect in some
cases~\cite{hinrichs2024blind}. They also found that SDR does not track
perceived dryness for distortion, as did Imort et al.\ and the EG-VAE authors.

\subsection{Dry signals from other sources}

The literature uses three other sources of a dry signal, each with a
mismatch. A \emph{different performance}, such as a corpus of DI guitar or the
player's own guitar, gives correct timbre statistics but different notes, so
it supports only unaligned losses, as in unpaired
modelling~\cite{wright2023adversarial}. A \emph{resynthesis} of the notes gives
the same notes at the same times but a different timbre: GuitarFlow turns
sample-based renders of tablature into realistic DI with flow
matching~\cite{loth2025guitarflow}, although for data generation rather than
tone estimation. We found no work that resynthesises a DI from transcribed
notes in order to fit a tone. A \emph{removal model} gives the same
performance with removal artefacts, as above.
